\documentclass[12pt]{article}
\usepackage{amsmath, amssymb}
\usepackage{geometry}
\usepackage{enumitem}
\geometry{margin=1in}

\title{Algebra II Final Practice: Conics, Statistics, Probability, and Functions}
\author{}
\date{}

\begin{document}
\maketitle

\section*{Instructions}
Complete in 60 minutes. Show clear work. Use exact forms for conic features unless a decimal approximation is explicitly requested. For statistics, use population standard deviation and round it to the nearest hundredth.

\section*{Reference Facts}
\begin{itemize}[leftmargin=1.5em]
  \item Circle:
  \[
  (x-h)^2+(y-k)^2=r^2
  \]
  \item Parabolas:
  \[
  (x-h)^2=4p(y-k),\qquad (y-k)^2=4p(x-h)
  \]
  \item Ellipses:
  \[
  \frac{(x-h)^2}{a^2}+\frac{(y-k)^2}{b^2}=1,\qquad c^2=a^2-b^2
  \]
  \item Hyperbolas:
  \[
  \frac{(x-h)^2}{a^2}-\frac{(y-k)^2}{b^2}=1,\qquad
  \frac{(y-k)^2}{a^2}-\frac{(x-h)^2}{b^2}=1,\qquad c^2=a^2+b^2
  \]
  \item Z-score:
  \[
  z=\frac{x-\mu}{\sigma},\qquad x=\mu+z\sigma
  \]
  \item Counting:
  \[
  {}_nC_r=\frac{n!}{r!(n-r)!}
  \]
\end{itemize}

\section*{Problems}
\begin{enumerate}[leftmargin=*, itemsep=1.15em, topsep=0.8em]
  \item Identify each conic as a circle, parabola, ellipse, or hyperbola. Do not rewrite in standard form.
  \[
  x^2+y^2-6x+8y-11=0
  \]
  \[
  y^2-4y-12x+16=0
  \]
  \[
  9x^2+4y^2-54x+16y+61=0
  \]
  \[
  16y^2-25x^2-96y-50x-281=0
  \]

  \item Rewrite the circle in standard form. Then identify its center and radius.
  \[
  x^2+y^2+10x-12y+25=0
  \]
  Then write the equation of the circle whose diameter has endpoints $(-2,7)$ and $(8,-1)$.

  \item Rewrite the parabola in standard form. Then state its orientation, vertex, focus, directrix, and axis of symmetry.
  \[
  x^2+8x-12y+4=0
  \]

  \item For the ellipse
  \[
  \frac{(x-2)^2}{25}+\frac{(y+1)^2}{9}=1,
  \]
  state the center, orientation, vertices, co-vertices, and foci.

  \item Write an equation of the ellipse with center $(1,-3)$, vertices $(1,5)$ and $(1,-11)$, and foci $(1,2)$ and $(1,-8)$.

  \item For the hyperbola
  \[
  \frac{(y-4)^2}{16}-\frac{(x+2)^2}{9}=1,
  \]
  state the center, orientation, vertices, co-vertices, foci, and asymptotes.

  \item Write an equation of the hyperbola with center $(3,-2)$, vertices $(-2,-2)$ and $(8,-2)$, and foci $(-4,-2)$ and $(10,-2)$. Then state its asymptotes.

  \item The data set below gives quiz scores.
  \[
  42,\ 45,\ 47,\ 47,\ 50,\ 52,\ 54,\ 57,\ 61,\ 65
  \]
  Find the mean, median, mode, range, population standard deviation, five-number summary, and IQR.

  \item A set of final exam scores is approximately normal with mean $76$ and standard deviation $8$.
  \begin{enumerate}[label=(\alph*)]
    \item Find the $z$-score for a score of $90$.
    \item Find the exam score with $z=-1.25$.
    \item Interpret the meaning of the $z$-score from part (a) in a complete sentence.
  \end{enumerate}

  \item A club has $8$ juniors and $6$ seniors. A committee of $4$ students is chosen at random.
  \begin{enumerate}[label=(\alph*)]
    \item How many committees are possible?
    \item How many committees have exactly $2$ juniors and $2$ seniors?
    \item What is the probability that the committee has exactly $2$ juniors and $2$ seniors?
  \end{enumerate}

  \item A parking permit code uses $2$ different uppercase letters followed by $3$ digits. The first digit cannot be $0$, and digits may repeat.
  \begin{enumerate}[label=(\alph*)]
    \item How many permit codes are possible?
    \item How many permit codes are possible if the digits may not repeat?
  \end{enumerate}

  \item A theater sells $300$ tickets when the price is $\$12$. For each $\$1$ increase in price, the theater expects to sell $15$ fewer tickets. Let $x$ be the number of $\$1$ price increases.
  \begin{enumerate}[label=(\alph*)]
    \item Create a price function $p(x)$ and a ticket-sales function $t(x)$.
    \item Create a revenue function $R(x)$.
    \item Find the ticket price that maximizes revenue.
    \item Find the maximum revenue and the number of tickets sold at that price.
    \item State a reasonable practical domain for $x$.
  \end{enumerate}
\end{enumerate}

\end{document}
